\documentclass[11pt]{article}

\usepackage[margin=1in]{geometry}
\usepackage{amsmath, amssymb, amsthm}
\usepackage{booktabs}
\usepackage{enumitem}
\usepackage{hyperref}
\usepackage{xcolor}

\title{\textbf{Project Plan and Paper Structure}\\
How Stable Are Regression Conclusions Across Subpopulations?\\
A Finite-Sample Empirical Study}
\author{Draft planning document}
\date{\today}

\begin{document}

\maketitle

\section{Working Title}
\textbf{How Stable Are Regression Conclusions Across Subpopulations? A Finite-Sample Empirical Study}

\section{Short Answer: Is a Literature Review Necessary?}
Yes, but only a \textbf{small and focused} one.

For this project, a long standalone literature review is probably unnecessary. A short review embedded in the introduction is enough if it does the following:
\begin{enumerate}[label=(\arabic*)]
    \item explains why pooled regression across heterogeneous groups can be misleading,
    \item briefly connects the paper to existing work on subgroup heterogeneity, model transportability, and distribution shift,
    \item identifies the gap your paper addresses,
    \item clarifies your empirical contribution.
\end{enumerate}

A good target is roughly \textbf{3--6 paragraphs} in the introduction devoted to prior work and positioning.

\section{Project Goal}
The goal of this paper is to study how stable the conclusions of a pooled linear regression are across subpopulations in finite samples, and when subgroup-specific departures make that pooled view break down.

More precisely, the paper asks:
\begin{quote}
When a linear regression model is fit on pooled data, how stable are the resulting conclusions---such as coefficient signs, coefficient magnitudes, and predictive performance---when checked against subgroup-specific fits and transported to other subgroups?
\end{quote}

The project will combine:
\begin{itemize}
    \item a controlled \textbf{simulation study}, and
    \item a small number of \textbf{real-data case studies}.
\end{itemize}

\section{Core Research Questions}
The paper will focus on the following questions:
\begin{enumerate}[label=\textbf{RQ\arabic*.}]
    \item Do coefficient \textbf{signs} remain stable across subpopulations?
    \item How much do coefficient \textbf{magnitudes} drift across subpopulations?
    \item When does a \textbf{pooled model} hide subgroup-specific relationships?
    \item To what extent can \textbf{interaction models} recover subgroup heterogeneity?
    \item How different is \textbf{predictive performance across groups} when a model is transported from one subgroup to another?
\end{enumerate}

\section{Main Working Hypothesis}
The main hypothesis is:
\begin{quote}
Even when a pooled linear regression appears adequate in aggregate, the regression conclusions may be unstable across subpopulations due to subgroup heterogeneity in coefficient values, predictor distributions, and noise structure.
\end{quote}

A secondary hypothesis is:
\begin{quote}
Predictive performance may remain moderately stable even when coefficient interpretation is not, so predictive adequacy alone does not guarantee interpretive transportability.
\end{quote}

\section{Main Contribution}
This paper is not intended to propose a new regression estimator. Instead, its contribution is primarily empirical and methodological.

The intended contributions are:
\begin{enumerate}[label=(\arabic*)]
    \item a finite-sample empirical study of how a pooled regression behaves across subpopulations,
    \item a comparison of a pooled baseline, subgroup-specific, and interaction-based models,
    \item a small set of interpretable stability metrics,
    \item practical guidance on when pooled regression conclusions should be treated with caution.
\end{enumerate}

\section{Proposed Stability Metrics}
The study will use a small number of simple and interpretable metrics.

\subsection{Sign Stability}
For each predictor, compare the sign of the estimated coefficient across subgroups and with the pooled estimate.

A simple summary metric can be:
\[
\text{Sign Stability Rate}_j
=
\frac{\#\{\text{comparisons where the sign agrees}\}}{\#\{\text{all relevant comparisons}\}}.
\]

\subsection{Standardized Coefficient Drift}
For predictor $j$, compare subgroup-specific estimates with the pooled estimate:
\[
D_j = \max_g \left| \hat{\beta}_{j,g} - \hat{\beta}_{j,\mathrm{pooled}} \right|.
\]

Report the standardized drift
\[
D_j^{\star} = \frac{D_j}{\widehat{\mathrm{SE}}\!\left(\hat{\beta}_{j,\mathrm{pooled}}\right)}
\]
to make comparisons across predictors more interpretable.

\subsection{Context-Normalized Drift}
To interpret drift relative to subgroup mean separation, define a predictor-level mean-spread index:
\[
M_j = \frac{1}{\binom{G}{2}} \sum_{a<b} \frac{\left|\bar{x}_{j,a} - \bar{x}_{j,b}\right|}{s_{j,\mathrm{pooled}}}.
\]
Then report context-normalized drift:
\[
\widetilde{D}_j = \frac{D_j^{\star}}{1 + M_j}.
\]
This helps distinguish large drift due to strong coefficient instability from large drift that is partly expected under strong subgroup mean separation.

\subsection{Transported Prediction Gap}
Train a model on one subgroup and evaluate it on another subgroup. Report metrics such as RMSE or MAE. The gap in out-of-group performance will quantify loss of transportability.

\subsection{Covariate Shift Summary}
Quantify predictor distribution differences across groups using a compact summary, such as the average standardized mean difference across predictors. This provides context for when cross-group transport failures are likely.

\section{Models to Compare}
The paper will compare the following models:

\subsection*{Model 1: Pooled OLS}
\[
Y = \beta_0 + X^\top \beta + \varepsilon.
\]

\subsection*{Model 2: Pooled OLS with Group Indicator}
\[
Y = \beta_0 + X^\top \beta + \gamma G + \varepsilon.
\]

\subsection*{Model 3: Pooled OLS with Group Interactions}
\[
Y = \beta_0 + X^\top \beta + \gamma G + X^\top \delta G + \varepsilon.
\]

\subsection*{Model 4: Subgroup-Specific OLS}
A separate linear regression will be fit within each subgroup.

Optional extensions such as ridge or lasso may be considered only if time permits, but they are not central to the project.

\section{Simulation Study}
The simulation study is the core of the paper.

\subsection{Basic Setup}
Generate data for $G$ subgroups, each with predictors and response
\[
X_g \sim \mathcal{N}(\mu_g, \Sigma_g), \qquad
Y_g = \beta_{0,g} + X_g^\top \beta_g + \varepsilon_g,
\]
where
\[
\varepsilon_g \sim \mathcal{N}(0, \sigma_g^2).
\]

\subsection{Simulation Scenarios}
The study will consider a small number of scenarios:

\begin{enumerate}[label=\textbf{S\arabic*.}]
    \item \textbf{No heterogeneity:} same coefficients, same predictor distribution, same noise across groups.
    \item \textbf{Covariate shift only:} same coefficients, different predictor distributions across groups.
    \item \textbf{Coefficient heterogeneity only:} different coefficients, same predictor distributions.
    \item \textbf{Combined heterogeneity:} different coefficients and different predictor distributions.
\end{enumerate}

If time allows, one additional scenario may be added with different noise levels across subgroups.

\subsection{Simulation Factors}
The simulation will vary a small set of factors:
\begin{itemize}
    \item sample size,
    \item number of predictors,
    \item degree of collinearity,
    \item signal-to-noise ratio,
    \item subgroup balance.
\end{itemize}

\subsection{Simulation Outputs}
For each replicate, the study will record:
\begin{itemize}
    \item estimated coefficients,
    \item coefficient signs,
    \item standard errors and confidence intervals,
    \item predictive performance within and across groups,
    \item interaction detection behavior.
\end{itemize}

\section{Real-Data Case Studies}
The empirical section should remain small and focused. Two or three real datasets are enough.

Desired dataset properties:
\begin{itemize}
    \item continuous response variable,
    \item at least one clear subgroup variable,
    \item multiple predictors,
    \item manageable missingness and preprocessing requirements,
    \item interpretable context.
\end{itemize}

Examples of subgroup variables include sex, age category, hospital/site, region, cohort, or socioeconomic category.

The purpose of the real-data section is not to establish universal truth, but to illustrate the practical relevance of the simulation findings.

\section{Why This Topic Matters}
This project matters for at least three reasons:
\begin{enumerate}[label=(\arabic*)]
    \item In applied work, analysts often fit a single pooled regression and interpret the coefficients as if they represented one stable relationship.
    \item In reality, subgroup heterogeneity may make such interpretations unreliable, even when the overall fit looks acceptable.
    \item In many practical settings, predictive performance alone is not enough; analysts also care about whether the sign and size of an effect are stable across contexts.
\end{enumerate}

\section{Mini Literature Review}
This review should be short and functional, not exhaustive.

\subsection{Suggested Themes to Cover}
The literature discussion in the introduction should briefly cover the following themes:

\begin{enumerate}[label=(\arabic*)]
    \item \textbf{Classical regression and heterogeneity.}  
    Standard linear regression often assumes a common linear relationship across observations unless interactions or stratification are introduced.

    \item \textbf{Subgroup effects and interaction modeling.}  
    It is well known that subgroup-specific relationships may be masked in pooled analyses, and interaction terms are a natural way to represent effect heterogeneity.

    \item \textbf{Transportability and distribution shift.}  
    Modern statistics and machine learning increasingly study how model conclusions and predictive performance change under population shift or domain shift.

    \item \textbf{Interpretation versus prediction.}  
    A model may predict reasonably well across groups while still producing unstable or misleading coefficient interpretations.

    \item \textbf{Gap addressed by this paper.}  
    Much of the existing literature is either broad and theoretical, or focused on general distribution shift. This paper instead gives a narrow finite-sample empirical study centered specifically on the stability of regression conclusions across subpopulations.
\end{enumerate}

\subsection{What the Literature Review Should \emph{Not} Try to Do}
Do not try to:
\begin{itemize}
    \item survey all of domain adaptation or all of distribution shift,
    \item review all heterogeneity modeling methods,
    \item make the paper look theoretical if it is mainly empirical,
    \item overclaim novelty by pretending no related work exists.
\end{itemize}

\section{Proposed Paper Structure}

\subsection*{1. Introduction}
This section should:
\begin{itemize}
    \item motivate the problem,
    \item explain why pooled regression may be misleading across subpopulations,
    \item provide a short literature review,
    \item identify the paper's gap,
    \item state the paper's contributions.
\end{itemize}

\subsection*{2. Methods}
This section should define:
\begin{itemize}
    \item notation,
    \item subgroup setting,
    \item models being compared,
    \item stability metrics,
    \item simulation design,
    \item evaluation protocol for real data.
\end{itemize}

\subsection*{3. Simulation Study}
This section should present:
\begin{itemize}
    \item scenario definitions,
    \item simulation setup,
    \item main figures and tables,
    \item interpretation of the observed instability patterns.
\end{itemize}

\subsection*{4. Real-Data Case Studies}
This section should:
\begin{itemize}
    \item describe the selected datasets,
    \item define subgroup variables,
    \item compare pooled and subgroup-aware models,
    \item connect the findings back to the simulation results.
\end{itemize}

\subsection*{5. Discussion}
This section should:
\begin{itemize}
    \item summarize the main findings,
    \item explain practical implications,
    \item discuss limitations,
    \item suggest future extensions.
\end{itemize}

\subsection*{6. Conclusion}
This should be short and direct.

\section{Practical Goals for the Project}
The project should aim to produce the following:

\begin{enumerate}[label=\textbf{G\arabic*.}]
    \item A reproducible simulation pipeline in R.
    \item A clear set of 3--5 main figures.
    \item A small number of interpretable summary tables.
    \item Two strong real-data illustrations.
    \item A paper draft that is concise, coherent, and clearly positioned.
\end{enumerate}

\section{What to Prioritize}
If time becomes limited, prioritize in this order:
\begin{enumerate}[label=(\arabic*)]
    \item simulation study,
    \item clear metrics,
    \item good figures,
    \item concise writing,
    \item real-data illustrations,
    \item optional extensions.
\end{enumerate}

\section{What to Avoid}
To keep the project publishable and finishable, avoid:
\begin{itemize}
    \item too many models,
    \item too many datasets,
    \item unnecessary deep learning baselines,
    \item excessive theory,
    \item broad claims about causality,
    \item trying to solve all of distribution shift.
\end{itemize}

\section{Minimal Viable Version of the Paper}
If the project must be compressed, the minimum viable paper is:
\begin{itemize}
    \item one clear introduction with a short literature review,
    \item one methods section,
    \item one simulation section with 4 scenarios,
    \item one real-data section with 2 datasets,
    \item one discussion/conclusion section.
\end{itemize}

\section{Final Recommendation}
Yes, include a literature review, but make it \textbf{short, focused, and embedded in the introduction}. For this paper, the strongest version is not one with the longest review, but one with:
\begin{itemize}
    \item a precise question,
    \item a clean simulation design,
    \item strong figures,
    \item a credible gap statement,
    \item and disciplined scope.
\end{itemize}

\end{document}